%% kapa: exam=1 round=22 year=2011 period=1 form=A subject=civil_law subject_ko=민법 seq=4 number=4
%% source: Q-Net 감정평가사 기출문제 (cst00309) artlSeq=5036790 file=extracted/1008/5036790/2011년 제22회 감정평가사 1차시험 기출문제/1교시.hwp
%% conversion: kordoc 4.24.51 -> tools/gen1.py (auto); review=proofread-sonnet; flags=none
\begin{question}{4}
2008년 3월 5일 탑승했던 비행기의 추락으로 생사불명이 된 甲에게 2010년 3월 5일 실종선고가 내려졌다. 乙은 甲의 배우자이다. 다음 설명으로 옳지 않은 것만을 모두 고르면?
\begin{framedtext}
ㄱ. 甲은 2008년 3월 5일 사망한 것으로 간주된다. \par
ㄴ. 2011년 4월 5일 甲이 종래의 주소지로 살아 돌아와 자신의 부동산을 丙에게 임대한 경우, 실종선고를 취소하지 않으면 그 임대차계약은 무효이다. \par
ㄷ. 2009년 4월 5일 선의인 乙이 丁에게 甲소유의 부동산을 양도한 경우, 甲의 생환을 이유로 실종선고가 취소되더라도 그 양도에는 영향을 미치지 않는다. \par
ㄹ. 乙이 甲의 토지를 상속하여 戊에게 2010년 4월 5일 매각한 후 甲이 살아 돌아와 실종선고가 취소된 경우, 乙과 戊가 선의라면 乙은 토지의 매매대금을 甲에게 반환할 필요는 없다.
\end{framedtext}
\begin{choices}
  \item ㄱ
  \item ㄴ, ㄷ
  \item ㄷ, ㄹ
  \item ㄱ, ㄴ, ㄹ
  \item ㄱ, ㄴ, ㄷ, ㄹ
\end{choices}
\end{question}
\answer{5}
